\chapter{Dipole-moment candidate batch (menu 19)}
\label{ch:menu19}
\menulabel{19}

\section{Purpose}

The scan behind the dipole-moment active-space selection: run the ground state
of every candidate space $\mathrm{A}_n = (n_\mathrm{e}, n_\mathrm{o})$ from one
cheap orbital source, compare each space's dipole moment with a reliable
reference, and keep the closest,

\[
\mathrm{dev}_n = \bigl| \mu(S_0, \mathrm{A}_n) - \mu(S_0, \mathrm{ref}) \bigr|,
\qquad \text{choose } \arg\min_n \mathrm{dev}_n ,
\]

a protocol whose numbers are in principle experimentally checkable.  This menu
writes the batch; menu~\menu{20} reads it and makes the choice.

\section{What you need}

The structure (XYZ) and a few settings.  The protocol is built for neutral
singlet ground states with a nonzero dipole moment; the menu refuses the rest
with the source's reason (the dipole of a charged molecule depends on the
coordinate origin).

\section{Walkthrough}

\lstinputlisting[style=fbkcode, firstline=54, lastline=65,
  title={The menu-19 example session (water: 35 candidates, MP2-orbital
  preparation)}]
  {examples/expected/m19_dm_batch.txt}

The batch's own files: the preparation input and the head of the manifest that
menu~\menu{20} will read (the candidate outputs are already named):

\lstinputlisting[style=fbkcode,
  title={examples/products excerpt: \file{h2o.fbk.dm/prep.inp}}]
  {examples/expected/products/h2o.fbk.dm/prep.inp}

\lstinputlisting[style=fbkcode, firstline=1, lastline=11,
  title={examples/products excerpt: the head of \file{h2o.fbk.dm/manifest.json}
  (35 candidates; the reference line is filled in after the runs)}]
  {examples/expected/products/h2o.fbk.dm/manifest.json}

\section{What you get}

A directory next to the structure (\file{<structure>.fbk.dm/}) with the
preparation input, one CASCI input per candidate, a run script and the
manifest.  The batch runs nothing -- this toolkit never does.

\section{How to read it}

\begin{readbox}
\begin{itemize}
  \item \textbf{The candidate family is the source's PASS set}: 6--14 active
    electrons (even), $n_\mathrm{e}/2 + 3 \le n_\mathrm{o} \le 14$ (35
    candidates); the PASS+ extras (the 4-electron and two-virtual rows) can be
    included -- the source kept them in its scanning set but flagged them as
    generally poor performers (51 candidates).
  \item \textbf{One orbital source for the whole scan}: the MP2 natural
    orbitals of the preparation run (\code{NatOrbs true}; they stay in the
    gbw the candidates read).  The source's own benchmark (its Figure~S5)
    found MP2 the best starting point for CASCI dipole moments; \code{hf} is
    the documented fallback.
  \item \textbf{Every candidate is a CAS-CI}: \code{!NoIter} with
    \code{moread} takes the input occupations and skips the orbital
    optimisation -- that is what makes the scan cheap (measured: ORCA prints
    ``MaxMacroIter 1 detected >>> CAS-CI''), and each single-root run prints
    its own dipole moment, which is exactly what the selection reads.
  \item \textbf{\code{nroots 1} is deliberate}: the protocol needs the
    $S_0$ dipole, and a state-averaged run prints only the state average.
\end{itemize}
\end{readbox}

\section{Boundaries and related sections}

\begin{refusalbox}
\begin{itemize}
  \item Charged systems and non-singlet ground states are refused (the
    source's own restrictions, with the reasons).
  \item The script is yours to run; the batch is text and nothing here
    touches an engine.
\end{itemize}
\end{refusalbox}

Related: \menu{20} (the selection that reads the batch), \menu{12} (the
entropy-based selection, a different criterion on the same question),
\menu{17} (carrying the chosen space across a geometry series), and
Chapter~\ref{ch:criteria}.
